% !PRINTED 204
\weberpage{204}

If $a$ is a $\nu$-fold pole of the group $G$ and

\weq{1}{a = \Theta(a) = \Theta^2(a)\cdots = \Theta^{\nu-1}(a),}

then by \S.\,52, 4. the second pole $a'$ of $\Theta$ is likewise $\nu$-fold,
and we have

\weq{2}{a' = \Theta(a') = \Theta^2(a')\cdots = \Theta^{\nu-1}(a').}

Now if the group $G$ has more than one system of conjugate
$\nu$-fold poles, then $a, a'$ may occur either in the same
or else in different ones of these systems.

But if there is only \gesperrt{one} system of $\nu$-fold poles, then $a$
and $a'$ must occur in the same system, and there must therefore
exist in $G$ a substitution $\psi$ not contained among the powers of $\Theta$,
such that

\weq{3}{a' = \psi(a)}

holds. From this it follows, by applying (1) and (2):
\[
\psi(a) = \Theta\psi(a), \quad a = \psi^{-1}\Theta\psi(a),
\]
from which we conclude that $\psi^{-1}\Theta\psi$ occurs among the powers of
$\Theta$, and hence, by (2), the following relations must hold:
\[
a' = \psi^{-1}\Theta\psi(a'), \quad \psi(a') = \Theta\psi(a')
\]
Thus $\psi(a')$ is also a pole of $\Theta$, and since $\psi(a')$
cannot be $= a'$, because otherwise $a'$ would be a pole of $\psi$, which
cannot be, since $\psi$ does not occur among the powers of $\Theta$,
we have

\weq{4}{\psi(a') = a.}

If now, by a transformation of the group, one brings the poles
$a$ and $a'$ to $0$ and $\infty$, then $\Theta$ takes the form
\[
\Theta(x) = \varepsilon x,
\]
in which $\varepsilon$ is a primitive $\nu$th root of unity, and $\psi$ must, on account of
(3) and (4), have the form:
\[
\psi(x) = \frac{c}{x},
\]
in which $c$ is a constant. By a further transformation
of the group one can give the constant $c$ any value whatever,
e.g. the value $1$, by transformation with the substitution $x\sqrt{c}$ for $x$.
We thus obtain the following theorem:

1.~\gesperrt{If in the group $G$ only one system of conjugate
 $\nu$-fold poles occurs, then one can}
%CONT>>
